\section{Confluent observation algebras and resolved directions}\label{sec:confluence}
The cardinality of a sumset is discontinuous when two of its elements
coalesce. The corresponding experiment need not be discontinuous. We now
identify the scales of the directions lost at such a collision. The
parameter $\theta$ below indexes a \emph{known calibration}; the unknown
parameter of each experiment is still $t$.

\subsection{Affine exponent families}
Fix $0=a_0<a_1<\cdots<a_{r-1}$ and real numbers $b_i$, with $b_0=0$.
On a sufficiently small compact interval $0\le\theta\le\theta_0\le1$ put
\[
 A_\theta=\{a_i+\theta b_i:0\le i<r\}.
\]
The order of these exponents is assumed unchanged. Take $I=[0,1]$ and
one fixed full-support prior $\mu$. A fixed coefficient matrix specifies
positive detector cells
\begin{equation}\label{eq:affine-calibration}
 k_j^\theta(t)=\sum_{i=0}^{r-1}c_{ji}t^{a_i+\theta b_i},
 \qquad \sum_jk_j^\theta=1.
\end{equation}
It has column rank $r$, and the cells have a common positive lower bound.
The comparator cube is fixed. These hypotheses include any sufficiently
small affine perturbation of a strictly positive calibration of full rank.
No continuity of an unknown prior density is assumed.

For a fixed positive future length $m\ge1$, let
\[
 \mathcal B_m=\left\{\left(\sum_{\nu=1}^m a_{i_\nu},
                              \sum_{\nu=1}^m b_{i_\nu}\right):
                              0\le i_\nu<r\right\}.
\]
Repeated \emph{pairs} are identified. Distinct pairs give distinct
exponents $\lambda+\theta\gamma$ for $0<\theta\le\theta_0$; shrink
$\theta_0$ to exclude crossings if necessary. For each first coordinate
$\lambda$, order its second coordinates as
$\gamma_{\lambda,0}<\cdots<\gamma_{\lambda,s_\lambda-1}$.
Set $K_m=\sum_\lambda s_\lambda$. The cluster at $\lambda=0$ consists
only of $(0,0)$.

Define the divided-difference functions
\begin{equation}\label{eq:divided-tests}
 \psi_{\lambda,k}^\theta(t)
  =[\lambda+\theta\gamma_{\lambda,0},\ldots,
       \lambda+\theta\gamma_{\lambda,k}]\,(x\mapsto t^x),
 \quad 0\le k<s_\lambda.
\end{equation}
At $\theta=0$ their continuous values are
\begin{equation}\label{eq:jet-tests}
 \psi_{\lambda,k}^0(t)=t^\lambda(\log t)^k/k!.
\end{equation}
For $\lambda>0$ these functions have value zero at $t=0$; the sole
zero-exponent function is the constant one. The bracket in
\eqref{eq:divided-tests} is a divided difference in the exponent, not a
new apparatus output. For a history $h$ let $z_\theta(h)$ be its posterior
expectations of these functions, omitting the constant coordinate.

List the orders $k$ of the nonconstant functions in nondecreasing order:
\[
 0=\nu_1\le\cdots\le\nu_d,\qquad d=K_m-1.
\]
There are $|mA_0|-1$ zero orders. The associated resolution function is
\begin{equation}\label{eq:resolution-function}
 \Phi_\theta(M)=\max_{1\le\ell\le d}
       \theta^{\,2(\nu_1+\cdots+\nu_\ell)/\ell}M^{-2/\ell},
       \qquad M\ge1.
\end{equation}
We set $0^0=1$ and $0^a=0$ for $a>0$ in this formula. In particular
$\Phi_0(M)=M^{-2/(|mA_0|-1)}$.

Choose $r$ attainable failure factors which form a basis in the formal
monomials $t^{a_i+\theta b_i}$, with coefficients independent of $\theta$.
For example use $1/2$ and $1/2+\delta t^{a_i+\theta b_i}$, $i>0$,
with $\delta$ sufficiently small. The fixed right inverse of the
calibration matrix realizes these by fixed command labels. All ordered
$m$-fold products form a common finite probe menu. Any other such fixed
basis, and its full product menu, gives equivalent bounds below.

\begin{theorem}[Uniform resolved checkpoint law]\label{thm:confluent-law}
Suppose
\begin{equation}\label{eq:full-future-condition}
 K_m-1\le n(r-1).
\end{equation}
Supply independent uniform commands for the first $n$ trials, then reveal
an independent uniform query from the preceding product menu. Let
$\overline R_{M;n,m}^\theta$ be the optimal unconditional squared-prediction
regret of an arbitrary $M$-message checkpoint encoder. Let
$R_{M;n,m}^{\theta,\max}$ be its minimax counterpart over admitted histories.
There are constants $0<c\le C<\infty$, independent of $M$ and $\theta$,
such that
\begin{equation}\label{eq:confluent-law}
 c\Phi_\theta(M)\le\overline R_{M;n,m}^\theta
       \le R_{M;n,m}^{\theta,\max}\le C\Phi_\theta(M),
 \qquad 0\le\theta\le\theta_0.
\end{equation}
The encoders may be Borel and randomized as in
Definition~\ref{def:finite-state}. At the common interior binomial command
tuple the normalized map $h\mapsto z_\theta(h)$ is a submersion of rank
$d$, including at $\theta=0$. Its local inverse bounds and an unconditional
minorization can be chosen uniformly in $\theta$.
\end{theorem}

The upper bound in this theorem is a checkpoint bound. It does not assume
that access to an exact prefix is free during repeated compression. Every
streaming filter is a checkpoint encoder, so the lower bound applies to
such filters as well. Section~\ref{sec:uniform} proves a matching genuinely
streaming bound across an entire singular family. Theorem~\ref{thm:attainable-checkpoint}
below removes the full-future restriction by truncating the resolution flag,
and Theorem~\ref{thm:general-uniform-streaming} proves the fixed-horizon
causal version for the affine class. When $m=0$ there is no nonconstant
future query and the prediction regret is zero; formula
\eqref{eq:resolution-function} is not used in that case.

\subsection{Confluent positivity and the physical prediction metric}
We give the details of the analytic and geometric steps in
Theorem~\ref{thm:confluent-law}.

\begin{lemma}[Strict confluent mixed pairing]\label{lem:confluent-positive}
Let $\alpha_1<\cdots<\alpha_s$ be nonnegative, and let $\lambda$ run
through distinct nonnegative numbers with multiplicities $q_\lambda$
summing to $s$. Include the columns
$t^\lambda(\log t)^k/k!$, $0\le k<q_\lambda$, in this order within
each increasing $\lambda$ cluster. Assume $q_0=1$ if zero occurs.
The matrix of their integrals against $t^{\alpha_i}\mu(\dd t)$ has
strictly positive determinant. Consequently the corresponding rectangular
pairing has full column rank whenever there are at least as many distinct
row exponents as columns.
\end{lemma}
\begin{proof}
Put $t=e^x$. The functions become $e^{\lambda x}x^k/k!$.
A nonzero linear combination with total multiplicity $s$ has at most
$s-1$ distinct zeros. Indeed divide by the first exponential and
differentiate $q_{\lambda_1}$ times. Rolle's theorem reduces a hypothetical
$s$ zeros to at least $s-q_{\lambda_1}$ zeros. The remaining terms are
exponential polynomials with distinct nonzero exponents and the same
remaining polynomial degrees: the operator $(D+\lambda-\lambda_1)^{q_{\lambda_1}}$
is invertible on polynomials of each fixed degree bound. Induction gives
a contradiction. If no remaining term occurs, the original function is
an exponential times a polynomial of degree less than $s$, and the
assertion is immediate.

Thus the evaluation determinant at distinct ordered points never
vanishes. Its sign is positive. One can determine the sign by letting
ordered simple exponents approach the prescribed repeated exponents and
dividing by their positive within-cluster Vandermonde factors. Equivalently,
the Wronskian of the displayed exponential-polynomial basis has sign
\[
 \prod_{\lambda<\lambda'}(\lambda'-\lambda)^{q_\lambda q_{\lambda'}}>0;
\]
the exponential factor and the factorial normalizations are positive.
Nonvanishing fixes this sign throughout the connected ordered-point domain.

Use the determinant integration identity from
Lemma~\ref{lem:mixed-moment}, with one determinant replaced by this
confluent evaluation determinant. Both determinants have the same positive
sign on ordered distinct positive points. The integrand is nonnegative
elsewhere by symmetry and continuity. Separated ordered subintervals have
positive prior mass, so the integral is strictly positive. The logarithmic
functions are bounded at zero under the stated hypotheses, which also
justifies Fubini. Select any $s$ rows for the rectangular assertion.
\end{proof}

\begin{lemma}[Desingularization and observable scales]\label{lem:observable-scales}
The functions in \eqref{eq:divided-tests}, and the needed derivatives in
$\theta$, are uniformly bounded on $[0,1]\times[0,\theta_0]$.
The posterior vectors $z_\theta(h)$ therefore lie in a fixed box.
For the common product probe menu let $p_\theta(h)$ be its probability
vector. There is a fixed matrix $G$ of full column rank such that
\begin{equation}\label{eq:physical-metric}
 p_\theta(h)-p_\theta(\widetilde h)
    =G\operatorname{diag}(\theta^{\nu_i})
          (z_\theta(h)-z_\theta(\widetilde h)).
\end{equation}
In particular the squared norm on the left is bounded above and below
by positive constants times
$\sum_i\theta^{2\nu_i}|z_i(h)-z_i(\widetilde h)|^2$,
with constants independent of $\theta$.
\end{lemma}
\begin{proof}
The divided difference of order $k$ is the integral of
$f^{(k)}(\sum_jw_jx_j)$ over the standard $k$-simplex, whose volume is
$1/k!$. This identity follows by iterating the fundamental theorem of
calculus, starting with $(f(y)-f(x))/(y-x)=\int_0^1f'(x+s(y-x))\dd s$.
It also gives the repeated-node value. Here $f(x)=t^x$, so every
integrand is $t^x(\log t)^k$. At a positive cluster the exponent is
bounded below by some $a>0$ uniformly on the fixed parameter interval.
For every fixed $j$, $\sup_{0<t\le1}t^a|\log t|^j<\infty$.
The same domination applies after finitely many $\theta$ derivatives.
There are no higher divided differences at zero. Posterior expectations
of bounded functions are bounded independently of the history.

Newton's interpolation identity, proved by successive subtraction of
the interpolating polynomials, is exact at each node:
\begin{equation}\label{eq:newton-exponents}
 t^{\lambda+\theta\gamma_{\lambda,j}}
   =\sum_{k=0}^j\theta^k
      \left(\prod_{v<k}(\gamma_{\lambda,j}-\gamma_{\lambda,v})\right)
                  \psi_{\lambda,k}^\theta(t).
\end{equation}
The matrix multiplying $\operatorname{diag}(\theta^k)$ is fixed,
triangular within each cluster, and has nonzero diagonal. The standard
identities used here are classical divided-difference calculus
\cite{deBoor}; the integral proof makes the uniformity at coalescence
explicit.

The expansion matrix from the formal pair monomials in $\mathcal B_m$
to the probe products has full column rank. To see this, invert the fixed
one-step basis change and multiply it $m$ times. Every formal monomial
product is a linear combination of the probe products. Identifying equal
pairs is a surjective linear quotient, so the spanning statement survives
that identification. Both this expansion matrix and the triangular
matrix in \eqref{eq:newton-exponents} are independent of $\theta$.
Their product has full column rank. Subtracting two normalized posterior
vectors removes its constant column and gives \eqref{eq:physical-metric}.
The remaining columns are still independent. Their extreme singular
values give the two constants. The equality extends to $\theta=0$ by
continuity; vanishing weights then remove precisely the unresolved jets.
\end{proof}

\begin{lemma}[Uniform attainable patch with its evidence]\label{lem:uniform-patch}
Under \eqref{eq:full-future-condition}, there are $\rho,\beta>0$,
independent of $\theta$, such that the subprobability distribution of
$z_\theta$ on all-failure prefixes dominates $\beta$ times uniform
probability on a translate of $[-\rho,\rho]^d$.
\end{lemma}
\begin{proof}
Use the binomial factors $(1+c_it^{D_\theta})/2$ of
Lemma~\ref{lem:binomial-tangent}, with distinct small $c_i$ and
$D_\theta=a_{r-1}+\theta b_{r-1}$. Since the coefficient calibration is
fixed, one common interior command tuple implements this family for every
$\theta$. Its unnormalized tangent has $n(r-1)+1$ distinct monomials,
also at zero. At positive $\theta$, pairing with the $K_m$ distinct
future exponents has rank $K_m$ by Lemma~\ref{lem:mixed-moment}.
The divided-difference basis change is invertible there. At zero,
Lemma~\ref{lem:confluent-positive} gives the same rank directly, without
passing a rank assertion through a singular limit. The product itself
lies in the tangent. The normalization map
$v\mapsto v-p e_0v$ has kernel exactly $\spanop\{p\}$ on this full
image, as in \eqref{eq:normalized-derivative}. Thus the normalized jet
map has rank $d$ at every $\theta$, including zero.

Here are uniform local details. Let $q$ be the total number of command
coordinates and $g_*$ the common interior tuple. The matrices
$J_\theta=D_gz_\theta(g_*)$ are continuous in $\theta$ and surjective.
Their least row singular value has a positive minimum $\sigma$ on the
compact parameter interval. Choose an orthonormal basis $E_\theta$ of
the row space and a complementary orthonormal basis $F_\theta$ of its
kernel, and use coordinates $g=g_*+E_\theta u+F_\theta w$.
The derivative at zero of
\[
 (u,w)\longmapsto
       (z_\theta(g_*+E_\theta u+F_\theta w)-z_\theta(g_*),w)
\]
is block diagonal with invertible upper block, least singular value at
least $\sigma$, and lower block the identity. Its derivative norm and
its second derivative have common finite bounds on a fixed neighborhood.
This follows from Lemma~\ref{lem:observable-scales}, finite product
rules, and the prefix evidence lower bound $\kappa^n$.

Choose a ball inside the command cube on which the derivative differs
from its central value by at most half its least singular value. This
radius can be chosen uniformly. For a sufficiently small target ball,
the map obtained by subtracting the nonlinear remainder and applying the
inverse central derivative is a contraction of that command ball into
itself. The contraction theorem gives a differentiable inverse. Hence a
product of a $d$-cube and a $(q-d)$-cube of common positive side length
lies in its image, uniformly in $\theta$. All inverse Jacobian
determinants are bounded below by a positive constant: the forward
derivative is uniformly bounded, so its determinant is uniformly bounded
above. The orthogonal command-coordinate change has determinant of
absolute value one. No continuity of the chosen orthonormal frames in
$\theta$ is needed for these uniform bounds.

Uniform commands have density $(1-2\eta)^{-q}$. Their joint law with the
all-failure word has density
\[
 (1-2\eta)^{-q}\mu\!\left(\prod_{i=1}^nF_{g_i}\right)
       \ge (1-2\eta)^{-q}\eta^n.
\]
Change variables using the inverse just constructed and integrate over
the complementary cube. This gives the asserted minorization, with
$\beta>0$ the product of this density bound, the inverse-Jacobian bound,
and the two cube volumes. The convention for a zero-dimensional
complement is volume one. This is a subprobability bound retaining the
failure evidence, not conditioning away a rare prefix.
\end{proof}

\subsection{A scale-sensitive finite-state calculation}
\begin{lemma}[Rectangular quantization at every budget]\label{lem:rectangle}
For $a_1\ge\cdots\ge a_d>0$ set
\[
 e_M(a)=\max_{1\le\ell\le d}
                  \left(\frac{a_1\cdots a_\ell}{M}\right)^{1/\ell}.
\]
A box with these side lengths has an at-most-$M$-cell rectangular
partition of diameter at most $2\sqrt d\,e_M(a)$ in every cell.
The optimal mean squared distortion of its uniform distribution using
$M$ arbitrary centers is at least $c_d e_M(a)^2$, with $c_d>0$
depending only on $d$. Zero side lengths may be omitted.
\end{lemma}
\begin{proof}
Put $e=e_M(a)$ and take $N_i=\max\{1,\lfloor a_i/e\rfloor\}$
equal intervals in coordinate $i$. If $\ell$ coordinates satisfy
$a_i\ge e$, then
$\prod_iN_i\le\prod_{i\le\ell}a_i/e^\ell\le M$.
Every side length in a cell is at most $2e$: for $a_i/e\ge1$ use
$\lfloor a_i/e\rfloor\ge a_i/(2e)$, and otherwise use $a_i<e$.
Choose an actual reachable representative in each nonempty cell when
covering a reachable subset of the box. The diameter bound remains valid.

For the lower bound project onto the first $\ell$ coordinates. The
projected uniform density is $(a_1\cdots a_\ell)^{-1}$ on its box.
The union of $M$ balls of radius $r$ has probability at most
$M v_\ell r^\ell/(a_1\cdots a_\ell)$, where $v_\ell$ is the volume of
the unit ball. Taking
$r=(a_1\cdots a_\ell/(2Mv_\ell))^{1/\ell}$ shows that mean squared
distance is at least $r^2/2$. Projection cannot increase distance.
Apply this for every $\ell$ and take the smallest of the finitely many
positive dimensional constants. Centers need not lie in the box.
\end{proof}

\begin{proof}[Proof of Theorem~\ref{thm:confluent-law}]
Let $D_\theta=\operatorname{diag}(\theta^{\nu_i})$; this temporary
matrix notation is unrelated to the peak dimension $D_A(N)$.
By Lemma~\ref{lem:observable-scales}, $D_\theta z_\theta(h)$ lies in
a translated box of side lengths at most $2B\theta^{\nu_i}$ for a fixed
$B$. Lemma~\ref{lem:rectangle} partitions this box into at most $M$
cells with squared diameter at most $C\Phi_\theta(M)$. Choose a
reachable representative in each nonempty cell. At the checkpoint,
encode the cell of the exact history and decode its representative's
probe probabilities. These are probabilities in $[0,1]$. Equation
\eqref{eq:physical-metric} gives the uniform upper regret bound.
Finite cells with deterministic boundary conventions give Borel encoders.

For the lower bound first average private decoder coins conditional on
the message and query. Squared loss cannot increase, leaving at most
$M$ centers in probe space. Orthogonally project these centers onto the
affine space $p_*+\operatorname{range}(G)$; projection cannot increase
distance to a true prediction. Invert $G$ on its range. Its lower singular
bound reduces the problem to distance from at most $M$ arbitrary centers
in the scaled $z$ coordinates. The law there dominates $\beta$ times
the image of the cube from Lemma~\ref{lem:uniform-patch}; its side lengths
are $2\rho\theta^{\nu_i}$. The lower half of
Lemma~\ref{lem:rectangle} gives $c\Phi_\theta(M)$ after division by
the fixed number of queries. Encoder randomization cannot beat distance
to the closest center at a fixed history. Fresh public randomness can be
fixed and then averaged because it is independent of the entire prefix.
At zero, omit zero-weight coordinates in this argument. The same constants,
reduced if necessary over the finitely many possible dimensions, work
throughout the compact parameter interval. This proves both criteria and
the claimed submersion statement.
\end{proof}

Thus the relevant finite-resolution data are the multiplicities of additive
collisions and their confluent orders, not only the cardinality of the
nearby sumset. The result is a uniform comparison up to constants. It
neither asserts a new general quantization principle nor identifies a
universal leading distortion constant.
